\par\addvspace{12pt}\Needspace{9\baselineskip}

\section{概率模型}

% ITEM: probability-example-007
\solutionitem{probability-example-007}{例6.7}
从编号1至5的球中依次不放回取两球。分别用有序对和无序集合表示样本空间，并判断哪种记录能表示事件“第一次取到1号球”。
\solutionanswer
有序空间为\(\{(i,j):1\le i,j\le5,i\ne j\}\)，有\(20\)点；无序空间为所有二元子集，有\(\binom52=10\)点。“第一次取到1”在前者为\(\{(1,j):j=2,3,4,5\}\)。后者不能表示这个事件，因为\((1,2)\)、\((2,1)\)被合成同一记录，却只有前者满足条件。


% ITEM: probability-example-008
\solutionitem{probability-example-008}{例6.8}
（2023年全国甲卷文科，题面整理）
某校文艺部有4名学生，高一、高二各2名。从中随机选2名组织文艺汇演，求两人来自不同年级的概率。
\solutionanswer
四个学生的二元子集等可能，共\(\binom42=6\)种；不同年级的选法为\(2\times2=4\)种，概率\(2/3\)。


% ITEM: probability-example-009
\solutionitem{probability-example-009}{例6.9}
（2024年全国甲卷文科，题面整理）
甲、乙、丙、丁各出场一次，次序由随机抽签确定。求丙不是第一个出场，且甲或乙最后出场的概率。
\solutionanswer
有利次序可先选最后出场的甲或乙，再在其余三人中选择非丙的第一人，最后排列剩余两人，共\(2\times2\times2=8\)种。全部\(4!=24\)种等可能，故概率\(1/3\)。


% ITEM: probability-exercise-03-01
\solutionitem{probability-exercise-03-01}{习题6.3（1）}
抛掷两枚公平骰子。分别用有序点数对与点数和记录结果，说明哪一种记录的样本点等可能，并求点数和为7的概率。
\solutionanswer
按两枚骰子的身份记录有序对，36个点在独立公平模型中等可能；点数和有11个可能值，却不等可能，例如和为2只对应一对、和为7对应六对。因此所求概率\(6/36=1/6\)。


% ITEM: probability-exercise-03-02
\solutionitem{probability-exercise-03-02}{习题6.3（2）}
设\(P(A)=0.6\)、\(P(B)=0.5\)、\(P(A\cup B)=0.7\)，求\(P(A\cap B)\)、\(P(A\setminus B)\)和\(P(A^c\cap B^c)\)。
\solutionanswer
加法公式得\(P(A\cap B)=0.6+0.5-0.7=0.4\)。继而\(P(A\setminus B)=0.2\)，\(P(A^c\cap B^c)=1-0.7=0.3\)。


% ITEM: probability-exercise-03-03
\solutionitem{probability-exercise-03-03}{习题6.3（3）}
从编号1至6的球中不放回抽取3个。分别按有序抽取和无序子集计数，求其中恰有两个偶数的概率。
\solutionanswer
无序计数得\(\binom32\binom31/\binom63=9/20\)。有序计数时分母\(A_6^3=120\)，有利子集各有\(3!=6\)个次序，分子为\(9\times6=54\)，仍得\(9/20\)。


% ITEM: probability-exercise-03-04
\solutionitem{probability-exercise-03-04}{习题6.3（4）}
某次模拟做1000次试验，事件发生397次。再做1000次时，是否必须恰好发生403次，才能使总频率等于模型给定的\(0.4\)？说明模型概率与实际记录各自表示什么。
\solutionanswer
若额外要求两批总频率恰为\(0.4\)，第二批次数确需\(2000\times0.4-397=403\)。但模型概率并不要求这项额外等式成立，也不强制第二批恰发生403次。\(0.4\)是机制给定的权重，\(397/1000\)以及后来的合计频率是随机记录形成的实际比例。
\par
